% Complete Experimental Table - Docker vs Apple container Benchmark
% Automatically generated by analyze_results_enhanced.py
%
% To use in LaTeX, copy the contents of the desired section
%
\documentclass{article}
\usepackage{booktabs}
\usepackage{longtable}
\usepackage{multirow}
\usepackage{array}
\begin{document}

% ============================================
% TABLE: Math
% ============================================

\begin{table}[h]
\centering
\small
\caption{Experimental Results: Math}
\label{tab:math}
\begin{tabular}{lccccccccc}
\toprule
Scenario & Tool & N & Mean & SD & Median & IQR & CI 95\% & Effect Size & p-value (FDR) \\
 & & & & & & & [lower, upper] & (Cohen's d) & \\
\midrule
\multirow{2}{*}{pi\_1t} & Docker & 30 & 84.0550 & 0.0483 & 84.0529 & 0.0274 & [84.0369, 84.0730] & 2.1907 & 0.0000*** \\
 & Container & 30 & 83.7033 & 0.2218 & 83.6689 & 0.0765 & [83.6205, 83.7861] &  &  \\
\midrule
\multirow{2}{*}{pi\_2t} & Docker & 30 & 87.4314 & 0.1717 & 87.3630 & 0.1770 & [87.3673, 87.4955] & 1.8351 & 0.0000*** \\
 & Container & 30 & 87.0556 & 0.2333 & 87.0178 & 0.0457 & [86.9684, 87.1427] &  &  \\
\midrule
\bottomrule
\end{tabular}
\begin{tablenotes}
\small
\item N: sample size; SD: standard deviation; IQR: interquartile range; CI: confidence interval; Effect Size: effect size (Cohen's d); p-value corrected by FDR (Benjamini-Hochberg).
\item Significance: *** p < 0.001, ** p < 0.01, * p < 0.05
\end{tablenotes}
\end{table}

\end{document}
